\documentclass[UTF8,11pt,a4paper,fontset=fandol]{ctexart}

\usepackage[left=2.4cm,right=2.4cm,top=2.5cm,bottom=2.5cm]{geometry}
\usepackage{amsmath,amssymb,amsthm,bm}
\usepackage{booktabs,longtable,tabularx,array,multirow}
\usepackage{graphicx}
\usepackage{tikz}
\usetikzlibrary{arrows.meta,positioning,fit}
\usepackage{xcolor}
\usepackage{enumitem}
\usepackage[numbers,sort&compress]{natbib}
\usepackage{hyperref}
\usepackage{microtype}
\usepackage{fancyhdr}
\usepackage{caption}
\usepackage{subcaption}
\usepackage{float}
\usepackage{url}

\definecolor{traceblue}{HTML}{2457A6}
\definecolor{tracecyan}{HTML}{DDECF8}
\definecolor{tracegreen}{HTML}{E5F2E8}
\definecolor{traceorange}{HTML}{FCE9D5}
\definecolor{tracegray}{HTML}{F3F5F7}
\hypersetup{
  colorlinks=true,
  linkcolor=traceblue,
  citecolor=traceblue,
  urlcolor=traceblue,
  pdftitle={从答案正确到结论可审计：面向 AI4Math 的 TRACE 分层可靠性框架},
  pdfauthor={匿名作者}
}
\setlength{\parindent}{2em}
\setlength{\parskip}{0.2em}
\setlist{nosep,leftmargin=2em}
\renewcommand{\arraystretch}{1.18}
\pagestyle{fancy}
\setlength{\headheight}{14pt}
\fancyhf{}
\fancyhead[L]{AI4Math 可靠性框架}
\fancyhead[R]{匿名投稿稿件}
\fancyfoot[C]{\thepage}
\captionsetup{font=small,labelfont=bf}

\newtheorem{proposition}{命题}
\newtheorem{definition}{定义}
\newcommand{\yes}{\textbf{Y}}
\newcommand{\partialyes}{\textbf{P}}
\newcommand{\no}{--}
\newcommand{\trace}{\textsf{TRACE}}

\title{\bfseries 从“答案正确”到“结论可审计”：\\面向 AI4Math 的 \trace{} 分层可靠性框架}
\author{匿名作者\\\small 独立研究者（匿名评审稿）}
\date{2026 年 9 月}

\begin{document}
\maketitle

\begin{abstract}
大语言模型在数学解题、定理证明与数学发现中表现出快速提升的能力，但“最终答案正确”不足以推出“推理过程可靠”：生成器可能偶然命中答案，学习型验证器可能与生成器共享偏差，形式化证明也可能证明了一个与原题不等价的命题，而公开分数可能因数据污染、预算不透明或环境漂移而无法复现。本文提出 \trace{}（\textbf{T}ask/Data Traceability、\textbf{R}easoning Search、semantic \textbf{A}lignment、\textbf{C}ertification、\textbf{E}valuation Replay）五门可靠性框架，将 AI4Math 系统的可信性拆解为任务与数据溯源、候选搜索与过程判别、语义规约对齐、确定性确证以及可复现实验五个连续门控。本文的核心观点是：形式化内核只能保证“给定形式化陈述下的证明有效”，不能自动保证自然语言问题与形式化陈述等价；因此，语义规约层是自然语言数学走向机器确证时的关键瓶颈。我们进一步给出门控式选择性发布规则、基于并集界的残余错误上界、风险--覆盖率与复现率等指标，并对 18 项代表性一手研究进行结构化证据编码。分析显示：搜索与验证能提高候选命中率，却不能把学习型评分器变成真值判定器；形式化证明提供最强的局部正确性保证，却把风险前移到陈述翻译、库版本和规格覆盖；公开基准改善可比较性，但仍需污染审计、预算披露和环境锁定。本文不报告新的模型性能，而贡献一个可证伪、可审计、可复用的系统评价框架及最小复现清单，为研究论文、模型评测和高风险数学应用提供共同的可靠性语言。
\end{abstract}

\noindent\textbf{关键词：} AI for Mathematics；数学推理；过程验证；自动形式化；定理证明；可复现评估；选择性预测

\begin{center}
\small\textbf{English Title:} From Correct Answers to Auditable Claims: The TRACE Layered Reliability Framework for AI4Math
\end{center}

\section{引言}
数学推理具有一个看似有利、实则容易被误解的特征：最终答案常常可被简短地判定。然而，答案级判定只观察输出末端，不能说明中间论证是否有效，更不能说明系统在题意理解、形式化翻译、工具调用与实验报告中是否保持了同一语义。以 GSM8K 为代表的工作展示了“多候选生成 + 验证器重排”的有效性\cite{cobbe2021verifiers}，MATH 则把评测推进到更复杂的竞赛题\cite{hendrycks2021math}；过程监督研究进一步表明，结果监督会把“过程错误但答案偶合正确”的轨迹视为正例\cite{uesato2022feedback,lightman2024verify}。另一方面，Lean 等证明助理提供小型可信内核，使形式化证明可以被机器确定性检查\cite{demoura2021lean4,mathlib2020}。这两条路线分别缓解“找不到正确候选”和“无法严格检查候选”的问题，却没有自动解决自然语言题意如何被无损映射到可检查对象的问题。

这一断点可概括为“验证对象缺口”：验证器验证的对象未必就是用户原本提出的命题。自然语言中的量词范围、隐含非退化条件、数域、单位、近似规则与图像约束，都可能在程序化或形式化时被省略或改变。自动形式化已显示出可行性\cite{wu2022autoformalization,azerbayev2023proofnet}，但其成功率和语义保真并非同一个指标。即使形式化证明完全通过，系统也可能只是严格证明了一个错误翻译后的命题。

本文研究三个问题：
\begin{enumerate}[label=\textbf{RQ\arabic*:}]
  \item 如何把“可靠”从单一正确率分解为可定位、可检验的系统属性？
  \item 如何连接概率式生成、学习型验证与形式化确证，同时避免把任一组件误当成全局正确性保证？
  \item 什么样的报告协议能够让第三方复查某项 AI4Math 结论的来源、语义、证明与实验环境？
\end{enumerate}

本文作出三项贡献。第一，提出 \trace{} 五门框架及跨层失效模型，将可靠性从模型属性改写为端到端证据链属性。第二，提出门控式选择性发布规则与残余错误并集界，明确系统应何时拒答，以及任何“可靠性分数”成立所需的假设。第三，对 18 项代表性一手研究进行机制编码，并给出可直接用于论文与基准报告的“可靠性卡”和最小复现清单。本文是一篇框架与证据综合论文，不声称进行了新的模型训练或性能对比。

\section{概念界定与研究方法}
\subsection{可靠性不是准确率的同义词}
给定自然语言任务 $x$、系统输出 $y$ 与目标语义 $\mathcal{S}(x)$，本文把可靠输出定义为：
\begin{equation}
\operatorname{Reliable}(x,y)=
\operatorname{Faithful}(x,z)\land
\operatorname{Valid}(z,\pi)\land
\operatorname{Grounded}(y;z,\pi)\land
\operatorname{Replayable}(m),
\label{eq:reliable}
\end{equation}
其中 $z$ 是结构化或形式化的中间表示，$\pi$ 是可检查证据（证明对象、程序运行轨迹、约束证书或逐步判定记录），$m$ 是包含数据、模型、预算、版本与随机性的实验清单。式~\eqref{eq:reliable} 区分四类常被混淆的事实：答案可能正确但过程无效；证明对象可能有效但陈述失真；自然语言解释可能与已验证对象不一致；一次运行可能成功但不可复现。

\subsection{文献选择与证据编码}
本文采用面向机制的结构化证据综合，而非声称穷尽全部 AI4Math 文献。纳入标准为：（1）直接提出数学生成、验证、形式化或评测方法；（2）有可核验的一手论文页、DOI 或 arXiv 标识；（3）能够回答至少一个研究问题。核心样本覆盖 2020--2026 年的 18 项研究，包括结果/过程验证\cite{cobbe2021verifiers,uesato2022feedback,lightman2024verify,wang2024mathshepherd}、数学语言模型与公开基准\cite{lewkowycz2022minerva,hendrycks2021math,he2024olympiadbench,glazer2024frontiermath}、形式化与自动形式化\cite{demoura2021lean4,mathlib2020,zheng2022minif2f,wu2022autoformalization,jiang2023draft,azerbayev2023proofnet,yang2023leandojo,poesia2023peano}，以及神经符号发现与证明系统\cite{trinh2024alphageometry,romeraparedes2024funsearch,hubert2026alphaproof}。

每项研究按 \trace{} 五门是否被\emph{显式实现或报告}编码为“覆盖”“部分覆盖”“未覆盖/不适用”。该编码评价的是证据链完备性，不是模型能力排名；“未覆盖”不等于工作质量低，而意味着该门不在原研究目标内。文献元数据以出版商、会议官方页面、Crossref 或 arXiv 记录交叉核验。为避免把综述性判断伪装成实验结果，本文不从图表反推未报告数字，不跨论文直接比较不同模型、提示、采样预算或数据切分下的准确率。

\section{相关工作}
\subsection{搜索与学习型验证}
生成器一次采样的失败可通过扩大测试时搜索得到缓解。给定候选集合 $C(x)=\{c_1,\ldots,c_k\}$，结果验证器通常选择
\begin{equation}
 c^*=\arg\max_{c\in C(x)} V_{\mathrm{out}}(x,c),
\end{equation}
其中 $V_{\mathrm{out}}$ 只依据最终答案或完整解答给分。该范式在数学文字题上建立了“生成--验证”基线\cite{cobbe2021verifiers}。但当错误推理偶然得到正确答案时，结果监督会提供错误的正反馈。过程奖励模型把轨迹 $c=(s_1,\ldots,s_T)$ 分解并评估中间步骤\cite{uesato2022feedback,lightman2024verify}；Math-Shepherd 则探索无需逐步人工标注的过程监督\cite{wang2024mathshepherd}。

这些方法的共同限制是：学习型验证器仍是统计模型。生成器与验证器可能继承相同训练分布、表述偏差或数据污染，增加采样预算还可能把搜索推向“最能欺骗评分器”而非最正确的候选。因此，学习型验证适合作为搜索启发式和风险筛选器，而不应被表述为形式化意义上的正确性证明。

\subsection{形式化证明与语义规约}
Lean 4 以小型内核检查证明对象\cite{demoura2021lean4}，mathlib 提供可复用的形式化数学基础设施\cite{mathlib2020}。miniF2F 使奥赛级形式化证明可以跨系统比较\cite{zheng2022minif2f}，LeanDojo 则把环境交互、检索、数据切分和开源工具链纳入可复现实验\cite{yang2023leandojo}。非正式证明草图可用于引导形式证明搜索\cite{jiang2023draft}，ProofNet 将问题扩展到本科数学的“自然语言到形式化再到证明”链路\cite{azerbayev2023proofnet}。

形式化路线提供的是\emph{相对于形式陈述和内核的可靠性}。自动形式化工作显示，大语言模型能够把部分自然语言题目转换为形式规范\cite{wu2022autoformalization}，但“语法可编译”“定理可证明”与“语义等价”并不相同。若形式化陈述漏掉约束，内核会忠实地确认错误规格。AlphaGeometry 采用专用几何语言与符号引擎，说明缩窄表示域能够提高可检查性\cite{trinh2024alphageometry}；其边界同样表明，表示覆盖率是可靠性的一部分。

\subsection{可复现评估}
MATH 提供分学科、分难度的竞赛题及解答\cite{hendrycks2021math}，OlympiadBench 增加双语、多模态和过程级评分需求\cite{he2024olympiadbench}，FrontierMath 通过保密题目降低污染风险，但集中式评测也限制了第三方本地复现\cite{glazer2024frontiermath}。这些工作体现了两种张力：公开数据便于复现却更易污染；保密数据有助于测量未见题能力，却降低审计透明度。因而，可靠评估不能只报告一个准确率，还应披露数据暴露风险、采样预算、判分器、工具版本和失败样本。

\section{\trace{}：五门分层可靠性框架}
\subsection{总体结构}
\trace{} 把端到端系统看作一条带拒答机制的证据流水线，如图~\ref{fig:trace} 所示。每一门都回答不同问题，任一门失败都应阻止“已验证”标签继续传播。

\begin{figure}[H]
\centering
\resizebox{\linewidth}{!}{%
\begin{tikzpicture}[
  node distance=8mm,
  box/.style={draw=traceblue,rounded corners=2mm,minimum height=16mm,text width=27mm,align=center,fill=tracecyan},
  gate/.style={draw=traceblue,circle,minimum size=7mm,fill=white,font=\bfseries},
  arr/.style={-{Latex[length=2.4mm]},thick,draw=traceblue}
]
\node[box] (t) {\textbf{T}\;任务/数据溯源\\题意、来源、污染审计};
\node[gate,right=3mm of t] (g1) {$g_T$};
\node[box,right=3mm of g1] (r) {\textbf{R}\;推理搜索\\多候选、过程判别};
\node[gate,right=3mm of r] (g2) {$g_R$};
\node[box,right=3mm of g2] (a) {\textbf{A}\;语义对齐\\Typed IR、双向核对};
\node[gate,right=3mm of a] (g3) {$g_A$};
\node[box,right=3mm of g3] (c) {\textbf{C}\;确定性确证\\内核、求解器、证书};
\node[gate,right=3mm of c] (g4) {$g_C$};
\node[box,right=3mm of g4] (e) {\textbf{E}\;评估重放\\版本、预算、日志};
\draw[arr] (t)--(g1); \draw[arr] (g1)--(r); \draw[arr] (r)--(g2);
\draw[arr] (g2)--(a); \draw[arr] (a)--(g3); \draw[arr] (g3)--(c);
\draw[arr] (c)--(g4); \draw[arr] (g4)--(e);
\node[draw=red!60,rounded corners,fill=traceorange,below=8mm of a,align=center,text width=56mm] (abstain) {任一门证据不足：降级为“未验证候选”或拒答\\不得输出“已证明/已复现”};
\draw[-{Latex[length=2mm]},dashed,red!70] (g1.south) |- (abstain.west);
\draw[-{Latex[length=2mm]},dashed,red!70] (g2.south) -- (abstain.north west);
\draw[-{Latex[length=2mm]},dashed,red!70] (g3.south) -- (abstain.north);
\draw[-{Latex[length=2mm]},dashed,red!70] (g4.south) -- (abstain.north east);
\end{tikzpicture}%
}
\caption{\trace{} 五门可靠性框架。圆形门控表示必须记录的验收判据，而非新的生成模块。}
\label{fig:trace}
\end{figure}

\subsection{T：任务与数据溯源门}
第一门固定“正在回答什么”。系统应保存原题、图像/附件、语言版本、时间戳与内容哈希；显式列出数域、量词、单位、非退化条件和允许使用的外部定理。评测时还需记录题目是否公开、是否可能进入预训练或微调数据，以及测试集是否在开发期间被查看。该门的主要失效是任务漂移与数据污染：系统对一个更容易的改写题给出正确答案，或在记忆题目上获得高分，却被误解为一般推理能力。

\subsection{R：推理搜索与过程判别门}
第二门负责提高“找到可验证候选”的概率，而非授予最终正确性。候选应保留采样温度、随机种子、最大 token、工具调用与搜索预算；过程验证器输出步骤分数和失败位置。对于 best-of-$k$，至少同时报告 $k$、生成成本和候选覆盖率，防止把计算预算差异解释为算法差异。过程监督相较结果监督提供更细粒度信号\cite{lightman2024verify}，但自动构造标签仍可能把生成器偏差带入验证器\cite{wang2024mathshepherd}。

\subsection{A：语义对齐与规约门}
第三门把自然语言 $x$ 和候选推理 $c^*$ 映射为类型化中间表示 $z=A(x,c^*)$。$z$ 可以是 Lean/Coq 陈述、符号程序、约束系统、抽象语法树或带类型的证明图。与只要求“能编译”不同，\trace{} 要求双向对齐：由 $z$ 反向生成规范化自然语言 $\hat{x}=D(z)$，并检查
\begin{equation}
\operatorname{Align}(x,z)=
\operatorname{Entails}(x,\hat{x})\land
\operatorname{Entails}(\hat{x},x)\land
\operatorname{ConstraintsEqual}(x,z).
\label{eq:alignment}
\end{equation}
其中蕴含检查可由领域规则、对抗性测试和人工审阅共同完成。人类签核必须面对“原题--规范化题意--形式陈述”三列对照，而不是只看可读性更高的自然语言解释。该门是全文强调的语义规约瓶颈：它位于概率生成与确定性验证之间，既不能由语言模型自我声明完成，也不能由证明内核代替。

\subsection{C：确定性确证门}
第四门要求使用与任务相匹配的独立检查器：形式证明由可信内核验证，代数恒等式由计算机代数系统与精确算术检查，组合构造由可执行判定器验证，数值结论则需要误差界或区间算术。FunSearch 将 LLM 的候选生成与可执行评价函数配对，展示了“难生成、易验证”问题上的有效闭环\cite{romeraparedes2024funsearch}；AlphaGeometry 通过神经辅助构造与符号演绎结合提供了更强的几何确证\cite{trinh2024alphageometry}；AlphaProof 则进一步展示了强化学习与 Lean 检查的结合\cite{hubert2026alphaproof}。但确证门的结论严格限定为：\emph{给定 $z$、库 $L$、内核 $K$ 与假设集 $H$，证据 $\pi$ 被接受}。它不自动覆盖 $x\leftrightarrow z$ 的语义等价、库公理选择或实现缺陷。

\subsection{E：评估重放门}
第五门要求第三方能够在隔离环境重放结论。最小清单包括：数据快照及哈希、训练/验证/测试切分、模型与权重版本、提示模板、解码参数、候选数与总算力预算、验证器版本、定理库与证明器版本、评分脚本、随机种子、硬件/软件环境、失败样本与人工裁决规则。LeanDojo 对环境、数据和模型的开源整合说明，形式化证明也需要工程层面的复现基础设施\cite{yang2023leandojo}。对于闭源模型或保密题库，应把“可重复获得相近统计结果”与“可逐项重放相同轨迹”分开报告。

\subsection{门控发布与残余风险上界}
定义发布变量 $G=\prod_{i\in\{T,R,A,C,E\}} g_i$，其中 $g_i\in\{0,1\}$ 表示第 $i$ 门证据是否满足预注册阈值。只有 $G=1$ 时，系统才能把结果标记为“经完整证据链验证”；否则应输出“未验证候选”或拒答。为了避免单一总分掩盖短板，建议报告可靠性向量
\begin{equation}
\bm{q}=(q_T,q_R,q_A,q_C,q_E),\qquad q_i\in[0,1],
\end{equation}
而不把它任意加权成一个看似精确的标量。

\begin{proposition}[门控链的保守错误界]
设 $F_i$ 表示第 $i$ 门发生未被发现的失效，并假设任何已发布的错误结论至少伴随一个 $F_i$。若在目标分布上 $\Pr(F_i)\leq\epsilon_i$，则
\begin{equation}
\Pr(\text{错误且发布})\leq \sum_i\epsilon_i.
\end{equation}
若覆盖率 $\kappa=\Pr(G=1)>0$，则
\begin{equation}
\Pr(\text{错误}\mid G=1)\leq \min\left(1,\frac{\sum_i\epsilon_i}{\kappa}\right).
\label{eq:unionbound}
\end{equation}
\end{proposition}
\begin{proof}
由前提，事件“错误且发布”包含于 $\bigcup_i F_i$。应用 Boole 不等式得到第一式；再除以 $\Pr(G=1)=\kappa$ 并用概率上界 1 截断，得到式~\eqref{eq:unionbound}。
\end{proof}

该命题的价值在于暴露假设，而非宣称现实系统已有可认证概率。若 $\epsilon_i$ 没有独立标注集、对抗测试或形式化保证支撑，就不能把式~\eqref{eq:unionbound} 当作数值证书。门控也必然牺牲覆盖率，因此可靠性评估必须同时报告风险与覆盖率。

\section{代表性研究的跨层证据分析}
表~\ref{tab:evidence} 对代表性工作进行机制级编码。符号 Y 表示论文明确实现或提供该门的主要证据，P 表示仅部分覆盖，-- 表示未覆盖或不适用。表格不是性能排行榜。

\begin{table}[H]
\centering
\scriptsize
\setlength{\tabcolsep}{3pt}
\caption{代表性研究对 \trace{} 五门的覆盖（基于原论文机制描述；Y=覆盖，P=部分覆盖，--=未覆盖或不适用）}
\label{tab:evidence}
\begin{tabularx}{\textwidth}{@{}>{\raggedright\arraybackslash}p{2.8cm}ccccc>{\raggedright\arraybackslash}X@{}}
\toprule
工作 & T & R & A & C & E & 主要残余风险 \\
\midrule
Verifier/GSM8K\cite{cobbe2021verifiers} & \partialyes & \yes & \no & \partialyes & \partialyes & 结果正确无法排除过程错误；公开题污染 \\
过程反馈\cite{uesato2022feedback} & \partialyes & \yes & \no & \partialyes & \partialyes & 步骤标签与任务外泛化仍有限 \\
PRM800K\cite{lightman2024verify} & \partialyes & \yes & \no & \partialyes & \partialyes & 学习型 PRM 不是逻辑判定器，标注成本高 \\
Math-Shepherd\cite{wang2024mathshepherd} & \partialyes & \yes & \no & \partialyes & \partialyes & 自动标签可能继承 rollout 模型偏差 \\
MATH\cite{hendrycks2021math} & \yes & \no & \no & \partialyes & \yes & 最终答案评分弱化过程可靠性；污染风险 \\
OlympiadBench\cite{he2024olympiadbench} & \yes & \no & \partialyes & \partialyes & \yes & 多模态与过程评分需严格复现裁决协议 \\
FrontierMath\cite{glazer2024frontiermath} & \yes & \no & \no & \partialyes & \partialyes & 保密题抗污染，但独立本地重放受限 \\
miniF2F\cite{zheng2022minif2f} & \yes & \partialyes & \partialyes & \yes & \yes & 形式陈述与自然语言题意的等价性需另审 \\
自动形式化\cite{wu2022autoformalization} & \partialyes & \partialyes & \yes & \yes & \partialyes & 可证明的翻译未必语义等价 \\
LeanDojo\cite{yang2023leandojo} & \yes & \yes & \partialyes & \yes & \yes & Lean/mathlib 版本与数据切分影响复现 \\
ProofNet\cite{azerbayev2023proofnet} & \yes & \partialyes & \yes & \yes & \partialyes & 本科题覆盖有限；需人工语义校对 \\
AlphaGeometry\cite{trinh2024alphageometry} & \partialyes & \yes & \yes & \yes & \partialyes & 专用表示覆盖范围窄，系统复现门槛高 \\
FunSearch\cite{romeraparedes2024funsearch} & \partialyes & \yes & \yes & \yes & \partialyes & 仅适用于存在强可执行评价函数的问题 \\
AlphaProof\cite{hubert2026alphaproof} & \partialyes & \yes & \yes & \yes & \partialyes & 形式化陈述、算力与闭源组件限制复现 \\
\bottomrule
\end{tabularx}
\end{table}

\subsection{发现一：搜索扩大“找到正确候选”的机会，不等于建立真值判定}
从结果验证器到过程奖励模型，研究重点逐渐从“答案是否匹配”转向“每一步是否可信”\cite{cobbe2021verifiers,lightman2024verify}。这是重要进步，但验证器与生成器同属学习系统时，二者错误并非独立：相同语料、相似架构与相同表述习惯会造成相关失效。因而，best-of-$k$ 的收益应与验证器的域外校准、对抗错误集表现和搜索预算同时报告。只报告最终准确率，无法区分“搜索覆盖提高”与“验证准确性提高”。

\subsection{发现二：形式化确证把风险从证明步骤前移到规格}
证明内核能排除非法推理步骤，却不能判断形式陈述是不是原题。自动形式化的进展\cite{wu2022autoformalization}和自然语言草图引导证明\cite{jiang2023draft}说明，语义规约不是一次性翻译，而是需要反向解释、约束测试与人工签核的独立环节。本文因此把 A 门与 C 门拆开：A 门问“证明的是不是同一件事”，C 门问“这个形式对象在指定系统内是否成立”。把二者合并会产生最危险的假阳性：机器给出完全通过的证明，但回答了错误问题。

\subsection{发现三：可复现性与抗污染性必须分别测量}
公开基准支持脚本复跑和横向比较，却可能被训练语料污染；保密基准减少记忆命中，却让题目、评分器与失败样本难以独立审计。两者不是简单的优劣关系。应分别报告（1）\textbf{重放透明度}：第三方能否在给定环境复现流程；（2）\textbf{新颖性可信度}：测试样本进入训练数据的风险；（3）\textbf{裁决独立性}：评分器是否独立于被评模型。把这三项压缩成“benchmark score”会隐藏重要差异。

\section{可复现实验与审计协议}
\subsection{最小实验设计}
若后续把 \trace{} 用于实证评估，建议采用分层、配对设计，而不是比较来源不同的排行榜数字。
\begin{enumerate}[label=\arabic*.]
  \item \textbf{任务分层：}至少包含答案型文字题、自然语言证明题、已形式化定理三类；按难度和领域分层抽样。
  \item \textbf{四个消融系统：}单次生成；best-of-$k$ + 结果验证；best-of-$k$ + 过程验证；完整 \trace{}（加入语义规约与确定性检查）。所有系统使用相同生成模型、提示和总 token 预算。
  \item \textbf{双盲语义审计：}两名审阅者独立标注形式陈述为“等价/偏强/偏弱/不相关”，分歧由第三名审阅者裁决，并报告一致性系数。
  \item \textbf{环境锁定：}保存容器镜像摘要、依赖锁文件、证明器与定理库提交号、评分脚本哈希、随机种子和失败日志。
  \item \textbf{预注册：}在查看测试结果前固定主要指标、拒答阈值、统计检验与排除规则；不得依据结果修改口径。
\end{enumerate}

\subsection{指标组}
可靠性应以向量而非单一准确率报告：
\begin{align}
\operatorname{AnsAcc} &= \frac{\#\text{最终答案正确}}{N},\\
\operatorname{ProcPass} &= \frac{\#\text{无已知无效步骤的轨迹}}{N},\\
\operatorname{SemEq} &= \frac{\#\text{经裁决语义等价的形式化陈述}}{\#\text{成功形式化}},\\
\operatorname{CertRate} &= \frac{\#\text{检查器接受且语义等价}}{N},\\
\operatorname{ReplayRate} &= \frac{\#\text{干净环境中成功重放}}{\#\text{重放尝试}}.
\end{align}
对于带拒答的系统，阈值 $\tau$ 下还应报告
\begin{equation}
\operatorname{Risk}(\tau)=\frac{\#\text{已接受但错误}}{\#\text{已接受}},\qquad
\operatorname{Coverage}(\tau)=\frac{\#\text{已接受}}{N}.
\end{equation}
只报告低风险而不报告覆盖率，会奖励“几乎全部拒答”的无用系统。成本维度至少包括每题候选数、总 token、工具调用次数、证明器 CPU 时间和人工审阅分钟数。

\subsection{可靠性卡}
每个公开结论应附一张机器可读或表格化的可靠性卡：
\begin{table}[H]
\centering
\small
\caption{\trace{} 可靠性卡的最小字段}
\begin{tabularx}{\textwidth}{@{}p{2.4cm}X X@{}}
\toprule
门 & 必填证据 & 失败时的发布标签 \\
\midrule
T & 原题/附件哈希、题意假设、数据来源、污染风险 & 任务来源不明 \\
R & 模型与提示版本、随机种子、$k$、预算、候选与过程分数 & 未验证候选 \\
A & Typed IR、反向释义、约束测试、人工语义签核 & 形式化未对齐 \\
C & 证明对象/程序证书、内核或求解器版本、检查日志 & 未确证 \\
E & 环境锁、评分脚本、完整配置、失败样本、重放记录 & 不可复现 \\
\bottomrule
\end{tabularx}
\end{table}

\section{讨论}
\subsection{框架的适用方式}
\trace{} 不是要求所有数学任务都立即形式化。对于小学算术，可执行程序和精确答案检查可能已足够；对于自然语言证明，应增加步骤级审阅与反例搜索；对于高风险定理声明，只有语义对齐后的形式证明或等价强度的独立证据才应获得“已确证”标签。关键不是一刀切地追求最重工具，而是让证据强度与结论风险相匹配。

该框架也解释了两类成功范式。Minerva 等模型说明领域数据与规模能显著提升定量推理能力\cite{lewkowycz2022minerva}，但能力提升本身不是可靠性证明。相反，AlphaGeometry、FunSearch 与 AlphaProof 的共同结构是把开放式生成限制在可被符号规则、执行器或证明内核检查的候选空间\cite{trinh2024alphageometry,romeraparedes2024funsearch,hubert2026alphaproof}。这类“生成负责提出，检查器负责否决”的分工比模型自我确认更稳健，但仍需 A 门与 E 门防止规格错位和不可复现。

\subsection{面向研究者与评测者的建议}
\begin{itemize}
  \item 把“证明通过”写成“在版本化形式陈述与指定内核下通过”，不要省略相对性。
  \item 把学习型验证器定位为风险模型，报告其校准、域外测试和与生成器共享数据的情况。
  \item 同时公布接受错误率与覆盖率；把拒答作为可靠系统的正常输出，而不是失败样本隐藏机制。
  \item 对自动形式化进行双向语义审计，尤其检查量词、数域、单位、边界与非退化条件。
  \item 把复现包视为论文论证的一部分，而非附录装饰；版本、预算和失败轨迹必须可追踪。
\end{itemize}

\section{局限性与效度威胁}
第一，本文是机制型框架和结构化证据综合，不是系统综述；文献样本具有目的性，可能遗漏其他证明助理、非英语工作和未公开工业系统。第二，表~\ref{tab:evidence} 的门覆盖编码只依据公开论文描述，不能证明未披露组件不存在；未来工作应由两名独立编码者计算一致性。第三，命题中的 $\epsilon_i$ 只有在目标分布、错误定义和估计协议固定时才有数值意义；现实中的跨层错误常相关，并集界可能很松。第四，形式化内核的可信性仍依赖实现、编译器、硬件和公理基础，不能被描述为形而上的“绝对正确”。第五，语义等价本身可能需要专家判断，难以完全自动化。第六，本文未运行新模型，因而不对任何系统的实际风险降低幅度作因果声明。

\section{结论}
AI4Math 的核心挑战不是让模型更频繁地产生看似正确的答案，而是让每个公开结论都携带一条可追溯、可检查、可重放的证据链。本文提出的 \trace{} 框架把可靠性分解为任务/数据溯源、推理搜索、语义对齐、确定性确证和评估重放五个门控，并用残余错误并集界说明端到端可靠性受最薄弱环节约束。跨研究分析表明，搜索与过程监督改善候选筛选，形式化内核强化局部正确性，公开基准支撑比较；但语义规约失真、验证器共偏差、数据污染和复现环境漂移仍不能由单一技术消除。未来工作应在同一预算下实证比较各门消融，并优先建立可审计的自然语言--形式语义对齐数据。对高风险数学主张，最诚实的系统不是从不拒答，而是清楚说明自己证明了什么、依据什么、在哪些条件下可以被再次验证。

\section*{AI 使用与研究诚信声明}
本文在资料整理、框架命名、结构草拟、LaTeX 排版与参考文献格式化中使用了生成式 AI。所有进入参考文献表的条目均通过出版商、会议官方页面、Crossref 或 arXiv 标识进行来源核对；未把模型生成的摘要当作一手证据。AI 曾把 DOI 含 2025 的 AlphaProof 论文误记为 2025 年卷期论文，经 Crossref 核对后按 Nature 第 651 卷第 8106 期的印刷年份 2026 记录，并保留 online/print 差异说明。论文不声称开展了未实际执行的模型实验；表格为机制编码，非性能复现。完整 AI 协作与核验轨迹见随稿的《AI协作日志.md》，七维复盘见《七维AAR.md》。投稿前仍应由作者本人完成逐条语义终审并签名确认。

\section*{数据与材料可用性}
本文不产生新的训练数据。论文源文件、BibTeX 数据库、文献核验表、AI 协作日志与 AAR 随稿提供。所有外部论文通过参考文献中的 DOI、官方会议页面或 arXiv URL 定位。

\bibliographystyle{unsrtnat}
\bibliography{references}

\appendix
\section{\trace{} 审计清单}
\begin{longtable}{@{}p{0.9cm}>{\raggedright\arraybackslash}p{5.0cm}>{\raggedright\arraybackslash}p{7.8cm}@{}}
\toprule
门 & 审计问题 & 最低通过证据 \\
\midrule
\endfirsthead
\toprule
门 & 审计问题 & 最低通过证据 \\
\midrule
\endhead
T & 原题是否完整、唯一、可追踪？ & 原始输入及哈希；量词、数域、单位、图像条件；数据来源与许可证 \\
T & 测试题是否可能被训练数据污染？ & 公开时间、去重方法、保密/动态题协议或污染风险声明 \\
R & 搜索是否公平比较？ & 同一模型与提示；相同 token/调用预算；$k$、温度、随机种子 \\
R & 验证器是否独立、校准且抗分布漂移？ & 独立测试集、校准曲线、对抗错误集、共享语料披露 \\
A & 形式化陈述是否与原题等价？ & 正反向释义；约束对照；边界/反例测试；人工签核 \\
A & 隐含条件是否被显式化？ & 非退化、定义域、单位、近似与容差清单 \\
C & 检查是否确定、可重放？ & 证明对象/证书；内核、求解器、库与版本；完整日志 \\
C & 检查器验证的结论边界是什么？ & 公理、依赖定理、数值误差或程序规格说明 \\
E & 第三方能否重建实验？ & 依赖锁、镜像摘要、配置、脚本、随机种子、硬件说明 \\
E & 负面结果是否保留？ & 全部失败样本、拒答、超时、人工裁决与排除原因 \\
\bottomrule
\end{longtable}

\section{建议的可靠性报告模板}
\begin{enumerate}
  \item \textbf{主张：}用一句可证伪的话描述，不使用“接近人类”等模糊措辞。
  \item \textbf{任务证据：}数据版本、样本范围、污染风险、授权与预处理。
  \item \textbf{生成与搜索：}模型、提示、解码、预算、候选保留规则。
  \item \textbf{语义规约：}IR/形式语言、双向对照、人工审阅者与分歧裁决。
  \item \textbf{确证：}检查器、版本、证书、失败类型和可信计算基。
  \item \textbf{评估：}主要/次要指标、置信区间、风险--覆盖率、成本。
  \item \textbf{复现：}环境、脚本、种子、日志、已知不确定性和可访问性。
\end{enumerate}

\end{document}
